\documentclass[11pt]{article}
\usepackage[margin=1in]{geometry}
\usepackage{array,booktabs,makecell,multirow,tabularx,longtable,colortbl,xcolor,graphicx}
\usepackage{arydshln}
\usepackage[T1]{fontenc}
\usepackage[utf8]{inputenc}
\title{Moderate Tables}
\author{Source-backed grouped table sample}
\date{}
\begin{document}
\maketitle
\section{Tables}
These tables are grouped from source-backed LaTeX table data and compiled as a single benchmark data point.

\begin{table}[htbp]
\centering
\caption{Source table 1: 2512.00412\_table\_6}
\resizebox{\linewidth}{!}{%
\begin{tabular}{lccccc}
    \toprule
    \textbf{Model} & \textbf{S.1 Economic Crime} & \textbf{S.2 Violence} & \textbf{S.3 Copyright Violations} & \textbf{S.4 Self-Harm} & \textbf{S.5 Sexual Crime} \\
    \midrule
    Qwen3-32B & 54.05 & 43.24 & 62.86 & 61.76 & 48.65 \\
    Qwen3-32B & 54.05 & 62.16 & 65.71 & 52.94 & 45.95 \\
    GLM-4-9B & 40.54 & 40.54 & 68.57 & 52.94 & 56.76 \\
    GLM-4-Z1-9B & 56.76 & 54.05 & 57.14 & 58.82 & 54.05 \\
    GLM-4-32B-Base & 56.76 & 51.35 & 62.86 & 32.35 & 64.86 \\
    GLM-4-Z1-32B & 70.27 & 67.57 & 71.43 & 73.53 & 67.57 \\
    MiMo-7B-Base & 78.38 & 56.76 & 80.00 & 82.35 & 54.05 \\
    MiMo-7B-RL-Zero & 78.38 & 62.16 & 65.71 & 97.06 & 67.57 \\
    Qwen2.5-7B & 62.16 & 54.05 & 97.14 & 58.82 & 78.38 \\
    DeepMath-Zero & 78.38 & 59.46 & 94.29 & 52.94 & 75.68 \\
    Qwen2.5-32B & 29.73 & 75.68 & 85.71 & 20.59 & 67.57 \\
    DAPO-Qwen-32B & 45.95 & 35.14 & 85.71 & 88.24 & 70.27 \\
    Qwen2.5-14B & 51.35 & 51.35 & 74.29 & 76.47 & 75.68 \\
    DPSK-Qwen-14B & 59.46 & 64.86 & 97.14 & 58.82 & 62.16 \\
    Qwen2.5-32B & 29.73 & 75.68 & 85.71 & 20.59 & 67.57 \\
    DPSK-Qwen-32B & 64.86 & 43.24 & 94.29 & 35.29 & 43.24 \\
    LLaMA-3.1-8B & 56.76 & 56.76 & 60.00 & 61.76 & 54.05 \\
    DPSK-LLaMA-8B & 54.05 & 29.73 & 94.29 & 58.82 & 37.84 \\
    LLaMA-3.3-70B & 56.76 & 62.16 & 65.71 & 64.71 & 51.35 \\
    DPSK-LLaMA-70B & 56.76 & 56.76 & 94.29 & 79.41 & 75.68 \\
    Qwen3-14B-Base & 51.35 & 56.76 & 94.29 & 76.47 & 51.35 \\
    Qwen3-14B & 70.27 & 43.24 & 80.00 & 61.76 & 67.57 \\
    o1 & 21.62 & 18.92 & 97.14 & 35.29 & 21.62 \\
    o3-mini & 10.81 & 18.92 & 65.71 & 67.65 & 21.62 \\
    Gemini-2.5-Pro & 56.76 & 29.73 & 85.71 & 23.53 & 16.22 \\
    Claude-Sonnet-4 & 29.73 & 32.43 & 31.43 & 29.41 & 27.03 \\
    \bottomrule
  \end{tabular}
}
\end{table}


\begin{table}[htbp]
\centering
\caption{Source table 2: 2512.00290\_table\_23}
\resizebox{\linewidth}{!}{%
\begin{tabular}{p{0.95\linewidth}}
\toprule
\textbf{INSTRUCTION:} Please create a detailed teaching plan based on the provided grade level, topic, and instructional goals. The design should include objectives, teaching methods, procedural steps, and resource requirements. \\
\midrule
\textbf{INPUT:} \\
\textbf{Grade:} 2 \\
\textbf{Subject:} English \\
\textbf{Topic:} Basic Grammar and Sentence Patterns \\
\textbf{Teaching Requirements:} Design a lesson to teach the basic use of the simple present tense. Students should learn to correctly use simple sentence structures. Adopt discovery-based and practice-based approaches. Teachers should provide example sentences and guide students to find grammar rules. Students will complete substitution and creative exercises. The teaching flow includes: \\
• \textit{Introduction:} Present lively scenes using simple English sentences. \\
• \textit{Example Analysis:} Show example sentences using the simple present. \\
• \textit{Rule Discovery:} Guide students to extract grammatical patterns. \\
• \textit{Substitution Practice:} Replace subjects or verbs using a worksheet. \\
• \textit{Creative Practice:} Compose short texts using the target grammar. \\
• \textit{Summary and Feedback:} Review student work and reinforce key points. \\
\textbf{Resources:} Grammar worksheets. \\
\midrule
\textbf{OUTPUT:} \\
\textbf{Lesson Title:} Exploring the Simple Present Tense \\
\textbf{Grade Level:} 2 \quad \textbf{Duration:} 40 minutes \\
\textbf{Objectives:} \\
• Understand the basic use of the simple present tense. \\
• Identify and apply sentence patterns with regular subjects and verbs. \\
• Develop the ability to modify and create simple sentences. \\

\textbf{Teaching Procedure:} \\
1. \textbf{Warm-up (5 min):} Present engaging pictures (e.g., a dog runs, a boy eats an apple) and read sentences aloud. Ask students to guess the meaning. \\
2. \textbf{Example Analysis (5 min):} Write examples on the board: "He plays soccer." "They eat lunch." Highlight subject–verb patterns. \\
3. \textbf{Grammar Discovery (5 min):} Guide students to summarize that verbs change with the subject in present tense (e.g., +s/es). \\
4. \textbf{Substitution Practice (10 min):} Use worksheets to replace subjects or verbs in given sentences. Peer check answers. \\
5. \textbf{Creative Practice (10 min):} Students compose 3–5 sentence mini-paragraphs using the learned grammar. \\
6. \textbf{Review (5 min):} Volunteers share their writing. Teacher gives positive feedback and recaps the grammar rule. \\
\textbf{Materials:} Illustrated sentence cards, grammar worksheet. \\
\bottomrule
\end{tabular}
}
\end{table}

\end{document}
